\begin{tikzpicture}[
  font=\sffamily\footnotesize,
  node distance=0.55cm and 0.45cm,
  etape/.style={
    draw=black!55,
    line width=0.7pt,
    rounded corners=2.5pt,
    align=center,
    minimum height=1.55cm,
    minimum width=2.45cm,
    inner xsep=2pt,
    inner ysep=5pt,
    fill=black!3,
    text=black!90
  },
  fleche/.style={
    -{Stealth[length=2.4mm, width=1.8mm]},
    line width=0.8pt,
    draw=black!50
  },
  badge/.style={
    circle,
    fill=black!78,
    text=white,
    font=\sffamily\bfseries\tiny,
    inner sep=0.6pt,
    minimum size=4.2mm
  }
]
  % Ligne supérieure
  \node[etape] (p1) {Document\\papier};
  \node[etape, right=of p1] (p2) {Capture\\numérique};
  \node[etape, right=of p2] (p3) {Extraction des\\informations};
  \node[etape, right=of p3] (p4) {Description};

  % Ligne inférieure
  \node[etape, below=1.2cm of p1] (p5) {Classement};
  \node[etape, right=of p5] (p6) {Conservation};
  \node[etape, right=of p6] (p7) {Recherche et\\exploitation};

  % Badges numérotés
  \foreach \n/\lab in {p1/1, p2/2, p3/3, p4/4, p5/5, p6/6, p7/7} {
    \node[badge] at (\n.north west) {\lab};
  }

  % Flèches horizontales
  \draw[fleche] (p1) -- (p2);
  \draw[fleche] (p2) -- (p3);
  \draw[fleche] (p3) -- (p4);
  \draw[fleche] (p5) -- (p6);
  \draw[fleche] (p6) -- (p7);

  % Passage de la ligne 1 à la ligne 2
  \draw[fleche] (p4.south) -- ++(0,-0.35) -| (p5.north);
\end{tikzpicture}
